\chapter{进一步的加热问题与已有求解尝试}\label{ch:further-heating}
前文已经说明，Hall场、物种电功、温升和不可逆熵产不能混为一谈。继续研究时，还应把问题从“离子怎样变热”推进到\textbf{两个物种如何相互制约、外部条件怎样改变能量路径、有限矩模型能保留哪些信息}。本章新增四篇主版本全文核读，共69个PDF物理页：Haggerty2015Partition、Dahlin2014Mechanisms、Dahlin2015ThreeD及原库已有PDF但此前未读的F627。它们均提供PIC/动理学参照，不能登记成四项已实现的Hall MHD加热闭合。F611、E553、F635等沿用已核读档案，不增计阅读篇数。

\begin{longtable}{p{.27\textwidth}p{.32\textwidth}p{.33\textwidth}}
\toprule 问题 & 已有求解尝试 & 目前还不能推出什么\\\midrule\endhead
电子与离子加热是否相互制约？ & 电子约束电势修正离子束增温；Haggerty等2015。 & 不能由测得电势的符合，推出无输入的普适分配闭合。\\
导引场如何改变电子增能渠道？ & 平行电场、Fermi与betatron分解；Dahlin等2014。 & 不能固定一种渠道，或把单粒子增能直接称作熵产。\\
三维改变总热量还是高能尾？ & 三维随机磁场增加加速区可达性；Dahlin等2015。 & 高能尾粒子增加十倍不等于电子热量增加十倍。\\
流剪切会不会使离子更热而喷流更慢？ & 上游流能、离子束选择与热反压；Haggerty等2025。 & 不能只用磁能归一，也不能凭弱喷流否定重联。\\
闭合能否预测加热对重联的反馈？ & 各向异性电子状态方程、十矩及热流闭合；F611等。 & 温度或总率单项接近，不证明跨参数能量分配正确。\\
多束混合、热核心与非热尾如何区分？ & 最大熵矩重建、多电子群十矩及相混合模型；E553、F635、E512。 & 高中心二阶矩、正分布或模型名称都不单独证明不可逆热化。\\
\bottomrule
\end{longtable}
表中证据分别展开如下；每节末的检验方案是本报告提出的研究设计，并非原文已经完成的工作。

\section{问题一：电子约束电势会不会同时限制离子加热？}
\textbf{为什么要问。}离子进入出流形成反向束，可以提高中心速度二阶矩；电子沿场快速运动，若要在出流中保持准中性并保留热电子，就需要相应的压力结构和平行电势。这个电势同时作用于离子。因此，先独立算一个离子拾取温升，再任意规定电子取得剩余能量，可能漏掉两物种之间的反馈。

\textbf{已有尝试。}\citet{Haggerty2015Partition}用对称、反平行2.5D PIC扫描上游磁场、密度和温度。图1中总温升在所扫参数内近似满足$\Delta(T_e+T_i)\simeq0.15m_iV_A^2$，分配却会改变；某些初始$T_e/T_i=9$的算例电子增温大于离子。这里$T$采用能量单位，$V_A$由上游重联磁场定义。\emph{0.15是温升对单粒子能量尺度的比例，不能直接写成15\%不可逆加热效率}；控制体热能、焓流和熵产还需分别计算。\citep{Haggerty2015Partition}

在忽略电子惯性、采用旋对称压力的排序下，沿场电子动量平衡给出
\begin{equation}
 eE_\parallel=-\nabla_\parallel T_{e\parallel}
 -T_{e\parallel}\nabla_\parallel\ln n
 +(T_{e\parallel}-T_{e\perp})\nabla_\parallel\ln B,
 \label{eq:question-epar}
\end{equation}
其中$e>0$。原文进一步把与密度梯度有关的电势差记为$\Delta\phi_n$，近似为
\begin{equation}
 e\Delta\phi_n\simeq T_{e\parallel,d}\ln(n_d/n_{\min}).
\end{equation}
它不是总平行电势的恒等式。原文根据几何和离子路径解释为何该部分对离子入流束重要，不能把所有电子温度梯度贡献都无条件归入同一个离子势垒。

在随收缩场线/出流运动的近似参考系里，令$v_0$为未受势垒减速的入射速度，$v_b$为越过有效势垒后的束速，则冷束估计为
\begin{equation}
 \tfrac12m_iv_b^2=\tfrac12m_iv_0^2-e\Delta\phi_n,
 \qquad
 \Delta T_i\simeq\frac{m_iv_b^2}{3}
 =\frac{m_iv_0^2-2e\Delta\phi_n}{3}.
 \label{eq:question-partition}
\end{equation}
温度的$1/3$来自三方向取迹，不表示已证明三方向均匀热化。此式只适用于冷入流束和相应轨道/势垒近似；右端为负时说明假设失效，不能预测负温升。图3显示未修正的拾取模型偏高，图4用模拟测得的电势修正后更接近PIC结果。电子被约束后可经历多次Fermi增能，电子留存与离子减速因而联系起来。

\textbf{解决到了哪里。}这是对物种分配的具体反馈机制，尚非闭合完成：图4输入了模拟电势，原文没有独立预测完整电子温升；简单慢激波估计的总温升约$0.25m_iV_A^2$，也不等于扫描中的约$0.15m_iV_A^2$。导场、三维、非对称边界和不可逆份额均未被这一扫描统一解决。其独立补充材料访问受限，完整参数表和部分轨道/平均方法未能另行核验。

\textbf{可检验的下一步。}在同一初边值下改变上游$T_e/T_i$与$\beta_e$，同时记录密度、$P_{e\parallel,\perp}$、沿场电势、离子束速及两物种焓流。先检查双温/十矩模型能否自洽产生电势及分配趋势，再检查电势修正是否只是事后拟合；把各向同性、各向异性和含热流模型分别比较，禁止逐个算例重新调耗散参数。

\loc{Haggerty2015Partition}{23页主版本全部；PDF 10--11式(1)--(4)，20--23图1--4；独立SI未取得}

\section{问题二：导引场怎样选择电子增能渠道与方向？}
\textbf{为什么要问。}同一个标量$T_e$可能来自平行加速、垂直压缩或非热尾增长。这些过程需要的压力和热流信息不同，反馈到出流电势与磁张力的方式也不同。一个只保留标量压强的Hall模型，即使温度峰值合适，也可能把正确数值建立在错误渠道上。

\textbf{已有尝试。}\citet{Dahlin2014Mechanisms}比较$B_g/B_0=0.2$与1的二维PIC，并用磁化电子的导引中心能量方程分析渠道。其非相对论形式可写为
\begin{equation}
 \frac{d\varepsilon}{dt}
 =q_eE_\parallel v_\parallel
 +\mu_e(\partial_t B+\vct u_E\cdot\nabla B)
 +m_ev_\parallel^2\vct u_E\cdot\vct\kappa,
 \label{eq:question-gc}
\end{equation}
其中$q_e=-e$，$\mu_e=m_ev_\perp^2/(2B)$，$\vct u_E=\vct E\times\vct B/B^2$（SI），$\vct\kappa=\vct b\cdot\nabla\vct b$。三项依次表示平行电场、磁场强度变化导致的betatron效应和曲率/Fermi效应。最后一项可理解为电子在收缩场线中反复反射而增能。这里追踪的是粒子动能，必须另作中心矩和分布诊断才能讨论热量；三项也不能再与完整$\vct J_e\cdot\vct E$相加，否则重复记账。

图5--6显示，弱导场时Fermi贡献主导，强导场时平行电场与Fermi贡献可比；磁场减弱时betatron可给净冷却。图9--11还提示，局地很强的$E_\parallel$经过空间积分可能部分抵消，岛合并使增能呈间歇性。图12的高能尾偏向平行方向，不能仅凭它得出电子已各向同性热化。\citep{Dahlin2014Mechanisms}

\textbf{解决到了哪里。}原文建立了可定位的渠道预算，但导引中心近似要求磁化电子，在早期小岛及合并的尖锐结构处会有偏差；局地贡献还受参考系影响。二维和质量比25的结果不能直接外推任意导场，也没有同时解决离子分配与熵增。

\textbf{可检验的下一步。}在同一导场序列下同时比较$T_{e\parallel}$、$T_{e\perp}$、热核心与尾部、场功和压力功。对十矩/各向异性模型检验各方向预算与场几何是否一致；若需要加上经验“Fermi加热项”，必须先确认原能量方程哪些贡献已经包含，避免重复回收能量。对于近零磁场区域另设近似失效诊断。

\loc{Dahlin2014Mechanisms}{19页主版本全部；PDF 11--12图5--6，15--17图9--11，18--19图12--13；图1--13全部核查}

\section{问题三：三维主要改变热核心，还是改变高能电子的逃逸与再加速？}
\textbf{为什么要问。}二维磁岛会把电子困在封闭磁面附近；三维磁场线连通性和粒子输运可以改变它们进入加速区的次数。能量供给相近并不保证能谱相同，局地热流闭合也未必记得粒子从哪里来、已经加速过多少次。

\textbf{已有尝试。}\citet{Dahlin2015ThreeD}比较含导场的二维、三维PIC。图2显示，在所比时刻二维释放的磁能反而约多15\%，但超过$0.5m_ec^2$的电子份额，三维约为$2\times10^{-4}$、二维约为$2\times10^{-5}$。因此\textbf{约十倍是指定阈值以上粒子数份额的差别，并非总电子热量的差别}。图3中三维高能电子分布范围更广；图4的轨道与所选最强750个电子支持更强输运。随机磁场让电子离开二维岛芯限制、接触更多增能区，高能电子的Fermi项更重要。\citep{Dahlin2015ThreeD}

\textbf{解决到了哪里。}这给出了三维影响高能尾的机制证据，尚未建立普适谱或逃逸定律。计算采用周期边界，没有真实系统的粒子逃逸；没有得到稳健幂律，电子也没有自动各向同性化。最强粒子轨道样本不能代表全部电子。主算例质量比25，虽补做其他尺寸和质量比100算例，仍不足以消除尺度限制；作者稿的两处初温写法不一致，详细笔记保留此版本问题，不据其中一个数值推导精确加热系数。

\textbf{可检验的下一步。}把热核心增温、超过固定能量阈值的粒子数/能量份额与总电子能量分开统计；同时固定输入功率，比较二维/三维、周期/开边界和多种系统长度。记录驻留时间与沿磁场热流。若只用一个$T_e$拟合出流，则需说明模型放弃了哪一类非热预测；若用非局地热流或多群模型，应检验它是否能同时复现尾部逃逸和总体能量守恒。

\loc{Dahlin2015ThreeD}{13页主版本全部；PDF 8--11图1--4；独立视频和补充图未取得}

\section{问题四：流剪切能否让离子更热、出流却更慢？}
\textbf{为什么要问。}存在上游沿场流剪切时，进入控制体的能量不止Poynting通量。若仍把全部下游热能除以释放磁能，就可能把流能转热误算成更高的磁能加热效率；若把离子Alfv\'en喷流当成唯一重联标志，又可能漏掉热反压显著的情形。

\textbf{已有尝试。}\citet{F627}用六组零导场、低$\beta_i=0.1$的2.5D PIC扫描$V_s/V_A=0$至1。图3所示离子$\Delta T_{i,xx}$约由$0.11$增到$0.35m_iV_A^2$，出流约由$0.62$降到$0.17V_A$。前者约为3.2倍，指特定方向的温度增量，不是标量温度、总离子热量或熵产同倍增长；最大名义剪切组的实际入流剪切也会衰减。原文文字称电子加热受影响不大，但未展示对应图，不能据此给出精确两物种加热比。

其解释在出流参考系中考察离子束选择、混合及散射：原有上游流动提供额外能量，较高离子热压又减慢出流。用于解释反压的近似动量关系是
\begin{equation}
 \frac{n_{\rm out}}{n_{\rm up}}
 \left(V_o^2+\frac{2\Delta T_{i,xx}}{m_i}\right)
 \simeq\epsilon V_A^2,\qquad
 \epsilon=1-\frac{\mu_0(P_\parallel-P_\perp)}{B^2}.
 \label{eq:question-shear}
\end{equation}
原文相关比较取$\epsilon\simeq1$、$n_{\rm out}/n_{\rm up}\simeq1.6$，还忽略了若干电子压强改变及小几何项；它不是独立的全能量守恒式。图5--6把理论与实际PIC扫描比较；其他$\beta$分支主要仍是理论预测。\citep{F627}

\textbf{解决到了哪里。}已有具体机制和参数序列支持“热得更多、喷得更慢”，但散射/各向同性化尚未被完整定量闭合。人工选择的出流平均范围及其标准差不是独立测量误差。论文约0.88的归一焓流指标，其分母未包含全部上游流能，不能翻译成88\%磁能不可逆转热。联系Parker Solar Probe弱喷流现象是待验证解释，不是本篇已完成的新观测确认。所读是2025年8月预印本，期刊身份已登记，正式版未逐页核读。

\textbf{可检验的下一步。}在双温Hall、十矩和PIC中使用相同流剪切与边界输入，分别列磁能、上游流能、两物种热流和出流功。比较$V_o$、$T_{i,xx}$、标量$T_i$及各向异性；压力反馈解释应同时符合动量和能量预算。观测识别要联合磁拓扑、电子响应和粒子分布，弱离子喷流或高温单独出现都不够。导场、三维、其他$\beta$和质量比收敛是后续必要检验。

\loc{F627}{14页主版本全部；PDF 4--5图1--2，7图3，9式(2)/图4，10与12图5--6}

\section{问题五：加热闭合能否正确预测对电势、层长和出流的反作用？}
\textbf{为什么要问。}加热不是仅在固定流场上附加一个温度输出。压力各向异性改变有效磁张力，电子压力决定部分平行电势，热流影响电子保留，离子热压又抵抗出流。模型可以在某一时刻给出相近重联率，却通过不同压力结构获得这个率。

\textbf{已有尝试。}已核读的\citet{F611}以动理学离子加各向异性电子闭合，对照等温电子与PIC；电子层长和面外磁场显著改变，而后期整体重联率仍可能接近。它没有独立完成两物种不可逆加热分配预测，二维磁通标签、俘获电子历史及Harris实现中的密度修正也限制直接移植。\citet{TenBarge2019MMS}演化十矩并输出物种热能，仍有各向异性位置和热流闭合误差；\citet{Finelli2021}说明流体电子闭合改变动理学离子响应，但将详细能量分配留待后续。这些是真正的闭合研究进展，证据强度与上述PIC机制参照不同。

\textbf{尚缺的组合检验。}需要同时预测电势、压力各向异性、出流、热流和物种预算，并跨上游温度、导场、剪切和尺寸保持有效。只调一个松弛尺度把总率或温度峰值拟合好，不能证明其余反馈链正确。数值耗散、超电阻能量回收和物种交换也必须明列，否则不同来源的“热”可能互相补偿。

\textbf{可检验的下一步。}先选一个PIC参照作参数标定，再冻结闭合参数，预测其他条件；逐项比较电势与分布矩，而非只比较最终温度。最小有用序列是：无剪切基准、提高$T_e/T_i$、加入导场、加入流剪切。三维和开边界随后增加，避免同时改变全部条件而无法定位失败原因。本建议是一项可实施的模型检验设计，尚无本文所核读论文完成全部组合。

\section{问题六：多束混合、热核心和非热尾需要多少流体自由度？}
\textbf{为什么要问。}两个各自很冷、平均速度相反的粒子群，合并后的中心二阶矩可以很大。有限分辨率又会抹平细束，造成近似Maxwell外观。模型必须说明自己在预测总中心矩、群内温度、非热尾，还是粗粒化/碰撞熵产；这些量的电子--离子比例不必相同。

\textbf{已有尝试。}\citet{E553}从PIC矩重建最大熵分布，比较5、10、14、21矩：更多矩改善平顶和三角主体，但21矩仍漏掉多束细节。它是分布重建检验，不能当作已经自洽演化21矩重联的证明；分布保持正也不等于证明物理熵产。\citet{F635}把电子分为两群、各演化十矩，改善部分热流结构，但仍有幅度误差，晚期每群又可能形成多束。其固定离子、分群方法和三维推广限制，使它尚不能给出两物种共同加热闭合。其“156倍”是存储量比较，不是已测运行速度倍数。

\citet{E512}及\citet{Loureiro2013Heating}提供另一种思路：保留高阶速度空间信息，追踪相混合输送的自由能，再由弱碰撞耗散；这对特定排序下的电子热化更有针对性，却不能不加扩展就解决完整离子竞争通道。\citet{Yoon2023PhaseSpace}对冷多束和混沌展宽的区分，则提供了检查温度定义的轨道参照。这几条路各自增加不同信息，还不能拼成现成通用模型。

\textbf{可检验的下一步。}用同一分布同时输出总温度、分群内温度、群间漂移能、尾部能量和可定义的熵诊断，再做空间/速度分辨率扫描。检查10矩、更多矩和多群模型在哪一步失效；加入能量守恒的弱碰撞后，检查群间漂移是否转成群内热能，以及耗散峰与电功峰是否分离。对于观测，应将仪器采样窗与分布卷积后再比较，不能把不同采样平均造成的束合并误认为模型已预测热化。

\section{优先推进顺序与尚待核读线索}
\textbf{第一优先：电子电势与离子反压。}先做问题一和四的统一控制体预算，能够直接检验“两物种独立给热源”的缺陷，也能避免遗漏上游流动能。\textbf{第二优先：导场与分布方向。}接着用问题二和六检查标量两温是否不足、十矩增加的自由度是否真能改善预测。\textbf{第三优先：三维输运。}在前两步预算可闭合后，用问题三检验开边界与驻留时间，最终以问题五的跨条件预测衡量闭合是否可迁移。这个顺序侧重可判定性，并不表示后面的物理不重要。

截至2026年11月13日的补充检索，还发现两条2026年线索：\emph{Role of ion acoustic instability in magnetic reconnection}（\href{https://doi.org/10.1017/S002237782510113X}{DOI: 10.1017/S002237782510113X}），可作为异常交换/微观不稳定性方向的待读候选；\emph{Quantifying Non-Maxwellian Properties of Ion Populations During Reconnection at Earth's Magnetopause}（\href{https://doi.org/10.1029/2026GL122882}{DOI: 10.1029/2026GL122882}），可作为问题六的观测诊断候选。\textbf{两篇尚未全文核读，只登记题名线索，不据摘要采纳其科学结论，也不计入144篇核读数。}Shay等2014年的电子温升参数研究（DOI: 10.1063/1.4904203）同样保留为前序待读连接点。

新增主版本、详细参数与限制见\path{questions_manifest.json}和逐篇JSON笔记；检索主题、下载与访问缺口见\path{research/ion_heating/questions_20261113/search_record.json}（仓库相对路径）。本章未运行原作者代码或重算原始轨道，提出的是有证据支撑的问题与后续检验方案。
